\documentclass[12pt]{article}
\usepackage[utf8]{inputenc}
\usepackage{amsmath, amssymb, amsfonts}
\usepackage{geometry}
\geometry{margin=1in}
\usepackage{enumitem}
\setlist[enumerate]{leftmargin=*}

\title{Quiz 1\\ \vspace{8pt} \Large ECON 308 -- International Economic Relations}
\author{Instructor: Muhebullah Karimzada}
\date{September 18, 2025}
\usepackage{comment}
\begin{document}
\maketitle


%==================================
\begin{enumerate}


\item Suppose there are two countries, \textbf{Home} and \textbf{Foreign}, and two goods, \textbf{Wine} and \textbf{Cheese}. 
The unit labor requirements are:

\begin{center}
\begin{tabular}{|c|c|c|}
\hline
Country & Wine (hours per unit) & Cheese (hours per unit) \\
\hline
Home & 2 & 1 \\
Foreign & 6 & 2 \\
\hline
\end{tabular}
\end{center}

Each country has 120 units of labor.

\begin{enumerate}[label=(\alph*)]
    \item Compute the \textbf{maximum production possibilities} (Wine and Cheese) for each country. 
    \item Calculate the \textbf{opportunity cost} of producing Wine in terms of Cheese for both countries. 
    \item Identify which country has a \textbf{comparative advantage} in Wine and which in Cheese. 
    \item If the world relative price of Wine is $P_W/P_C = 1.5$, determine which good each country will export.
\end{enumerate}

\bigskip

%==================================
\begin{comment}
\item Two countries, \textbf{Alpha} and \textbf{Beta}, produce \textbf{Shirts} and \textbf{Shoes} with the following unit labor requirements:

\begin{center}
\begin{tabular}{|c|c|c|}
\hline
Country & Shirts & Shoes \\
\hline
Alpha & 5 & 10 \\
Beta & 20 & 25 \\
\hline
\end{tabular}
\end{center}

Each country has 1,000 units of labor.

\begin{enumerate}[label=(\alph*)]
    \item Calculate the \textbf{autarky relative price} of Shirts in terms of Shoes in each country. 
    \item Suppose free trade sets the world relative price at $P_{Shirts}/P_{Shoes} = 0.6$. Determine the \textbf{pattern of trade}. 
    \item If Alpha specializes completely according to comparative advantage, compute the \textbf{total world output} of Shirts and Shoes before and after trade. 
    \item Show that both countries can \textbf{consume beyond their PPFs} after trade.
\end{enumerate}

\bigskip
\end{comment}
%==================================

\item Suppose trade flows follow the basic gravity equation:
\[
T_{ij} = A \cdot \frac{Y_i Y_j}{D_{ij}}
\]
where $T_{ij}$ is trade between country $i$ and $j$, $Y$ denotes GDP, $D$ is the distance between countries, and $A=1$. 

\bigskip
The data are as follows:

\begin{itemize}
    \item Country A GDP: \$2{,}000 billion
    \item Country B GDP: \$1{,}500 billion
    \item Country C GDP: \$500 billion
    \item Distance between A and B: 1,000 km
    \item Distance between A and C: 5,000 km
\end{itemize}

\begin{enumerate}[label=(\alph*)]
    \item Calculate the predicted trade flow $T_{AB}$ and $T_{AC}$. 
    \item Find the ratio of trade between A--B versus A--C. 
    \item Suppose a free trade agreement between A and C effectively \textbf{halves the distance cost} (so effective distance = 2,500 km). Recalculate $T_{AC}$ and explain how the agreement changes trade flows. 
\end{enumerate}
\end{enumerate}
\end{document}
